\documentclass[12pt,a4paper]{article}
\usepackage{fontspec}
\setmainfont{DejaVu Serif}
\setsansfont{DejaVu Sans}
\setmonofont{DejaVu Sans Mono}[Scale=0.85]
\usepackage[margin=2.2cm]{geometry}
\usepackage{amsmath}
\usepackage{amssymb}
\usepackage{xcolor}
\usepackage{enumitem}
\usepackage{booktabs}
\usepackage{array}
\usepackage{listings}
\usepackage{hyperref}
\usepackage{titlesec}

\titleformat{\section}{\large\bfseries}{}{0pt}{}
\setlength{\parskip}{0.5em}
\setlength{\parindent}{0pt}
\sloppy
\setlength{\emergencystretch}{3em}
\hypersetup{hidelinks}

\lstset{
  basicstyle=\ttfamily\small,
  keywordstyle=\color{blue!70!black}\bfseries,
  commentstyle=\color{green!45!black},
  frame=single, framerule=0.3pt, rulecolor=\color{gray!60},
  backgroundcolor=\color{gray!6},
  columns=fullflexible, keepspaces=true,
  language=C, showstringspaces=false,
  xleftmargin=1em, xrightmargin=1em
}

\newcommand{\question}[2]{\section{Дасгал #1}\textbf{Асуулт.} #2\par\medskip\textbf{Хариулт.}\par}

\begin{document}

\begin{center}
  {\Large\bfseries Лабораторийн ажил №10}\\[0.4em]
  {\large Их хэмжээний санах төхөөрөмж (Mass-Storage Structure) — Бүлэг 11, Practice Exercises}\\[0.8em]
  F.CSM302 Үйлдлийн систем\\
  Гүйцэтгэсэн: Б.Билгүүн (B190910014)\\[0.3em]
  Сурах бичиг: A. Silberschatz, P. B. Galvin, G. Gagne, \textit{Operating System Concepts}, 10th ed., Wiley, 2018
\end{center}
\hrule
\bigskip

Хатуу диск (HDD)-нд нэг хүсэлтийн хугацаа $\approx$ \textbf{хайлтын хугацаа} (seek — толгойг цилиндр рүү зөөх) $+$ \textbf{эргэлтийн хүлээлт} (rotational latency) $+$ дамжуулах хугацаа. Хайлт хамгийн их зардалтай тул \textbf{дискний хуваарилалт} (FCFS, SSTF, SCAN, C-SCAN) толгойн нийт шилжилтийг багасгахыг зорьдог. NVM (SSD) төхөөрөмжид хөдөлгөөнт хэсэг байхгүй тул ихэвчлэн FCFS (NOOP) хангалттай. Тооцоог \texttt{lab10/diskcalc.py}, хуваарилалтыг \texttt{lab10/disk\_sched.c}-ээр шалгасан.

%---------------------------------------------------------------
\question{11.1}{Нэг хэрэглэгчийн орчинд FCFS-ээс өөр дискний хуваарилалт хэрэгтэй юу? Тайлбарла.}

Нэг хэрэглэгчийн орчинд I/O дараалал ихэнхдээ хоосон эсвэл нэг хүсэлттэй байдаг — дараалалд сонгох зүйл байхгүй үед ямар ч алгоритм FCFS-тэй ижил ажиллана. Тиймээс ашиг нь бага.

Гэхдээ нэг хэрэглэгч ч \emph{олон процесс, урсгал} зэрэг ажиллуулдаг (браузер, файл хуулах, хайлтын индекс, вирусын шалгалт), мөн үйлдлийн систем өөрөө урьдчилан унших (read-ahead), хуудаслалт (swap), бичих буфер цэвэрлэх зэрэг олон хүсэлт үүсгэдэг. Ийм үед дараалалд олон хүсэлт хуримтлагдаж, SCAN/C-SCAN толгойн шилжилтийг мэдэгдэхүйц багасгана. Иймд ачаалал өндөр үед хэрэгтэй, ердийн үед ач холбогдол бага.

%---------------------------------------------------------------
\question{11.2}{SSTF хуваарилалт яагаад голын цилиндрүүдийг хамгийн дотор, хамгийн гадна цилиндрээс илүүд үздэгийг тайлбарла.}

SSTF толгойд хамгийн \emph{ойр} хүсэлтийг сонгоно. Дискний голд байгаа цилиндрээс бусад бүх цилиндр хүртэлх дундаж зай хамгийн бага (хоёр талдаа хүсэлттэй), харин зах дахь цилиндр зөвхөн нэг талдаа хүсэлттэй, дундаж зай нь хоёр дахин их. Тэгэхээр толгой голын бүсэд байх магадлал өндөр, тэнд шинээр ирэх хүсэлтүүд ойр байдаг тул тасралтгүй сонгогдоно; зах дахь хүсэлтүүд үргэлж “хол” болж, ачаалал ихтэй үед удаан хүлээж, хомсдолд (starvation) орж болно.

%---------------------------------------------------------------
\question{11.3}{Дискний хуваарилалтад эргэлтийн хүлээлтийг яагаад ихэвчлэн тооцдоггүй вэ? SSTF, SCAN, C-SCAN-ийг эргэлтийн хүлээлтийг оновчлох байдлаар хэрхэн өөрчлөх вэ?}

\textbf{Яагаад тооцдоггүй:} орчин үеийн дискүүд логик блокын хаяг (LBA) ашигладаг тул үйлдлийн систем блокын жинхэнэ физик байрлал (цилиндр, сектор), толгойн одоогийн эргэлтийн байрлалыг мэддэггүй; секторыг дахин буулгах (sparing), бүсчилсэн бичлэг зэргээс болж геометр таамаглах боломжгүй. Нөгөө талаас диск өөрөө (firmware) дотооддоо хүсэлтүүдийг дахин эрэмбэлж эргэлтийн хүлээлтийг оновчилдог. Мөн эргэлтийн хүлээлт (хэдэн мс) хайлтын хугацаанаас бага.

\textbf{Өөрчлөх арга} (эргэлтийн байрлалыг мэддэг гэж үзвэл):
\begin{itemize}[leftmargin=1.5em]
  \item \textbf{SSTF → SPTF (shortest positioning time first):} зөвхөн хайлтын зайгаар биш, \emph{хайлт + эргэлтийн хүлээлтийн} нийлбэр хамгийн бага хүсэлтийг сонгоно.
  \item \textbf{SCAN, C-SCAN:} цилиндрээр явах дарааллыг хэвээр үлдээж, нэг цилиндр дээрх олон хүсэлтийг (мөн ойрын цилиндрүүдийн хүсэлтийг) толгойн доор \emph{эргэлтээр ирэх} дарааллаар нь үйлчилнэ; сектор дөнгөж өнгөрсөн хүсэлтийг бүтэн эргэлт хүлээхийн оронд дараагийн цилиндрийн хүсэлтийг хийж байгаад эргэж ирж болно.
\end{itemize}

%---------------------------------------------------------------
\question{11.4}{Олон даалгаварт орчинд файлын системийн I/O-г дискнүүд ба контроллеруудын хооронд тэнцвэржүүлэх нь яагаад чухал вэ?}

Систем хамгийн удаан (хамгийн ачаалалтай) төхөөрөмжийн хурдаар хязгаарлагдана. Бүх идэвхтэй файл нэг диск/контроллер дээр байвал тэнд урт дараалал үүсч, хүсэлт бүрийн хүлээлтийн хугацаа өснө, бусад диск сул зогсоно. Ачааллыг тараавал дискнүүд \textbf{зэрэгцээ} ажиллаж, нийт дамжуулах чадвар (throughput) нэмэгдэж, хариу өгөх хугацаа багасна; нэг контроллерийн шугамын (bus) хурдыг ч хэтрүүлэхгүй. Энэ нь RAID-д өгөгдлийг хуваах (striping)-ын санаа юм.

%---------------------------------------------------------------
\question{11.5}{Кодын хуудсыг swap зайд хадгалахын оронд файлын системээс дахин уншихын давуу, сул тал юу вэ?}

\textbf{Файлын системээс дахин унших:} кодын хуудас зөвхөн уншигддаг, өөрчлөгддөггүй тул санах ойгоос гаргахдаа дискэнд \emph{бичих шаардлагагүй} — зүгээр хаяна; swap зай хэмнэнэ, гаргах ажил хурдан. \textbf{Сул тал:} файлын системээр унших нь удаан — файл дискэнд хуваагдмал байж болно, директор/inode хайлт, файлын системийн давхаргын ачаалал нэмэгдэнэ; мөн программ ажиллаж байхад файл өөрчлөгдөх/устахаас хамгаалах шаардлагатай.

\textbf{Swap ашиглах:} swap зай нь файлын системгүй, үргэлжилсэн блокоор, том хэсгээр зохион байгуулагдсан тул унших хурдан. \textbf{Сул тал:} эхний удаа кодыг swap руу \emph{бичих} нэмэлт I/O, swap зай их зарцуулна.

%---------------------------------------------------------------
\question{11.6}{Жинхэнэ тогтвортой санах ой (stable storage)-г хэрэгжүүлэх ямар нэг арга бий юу? Тайлбарла.}

\textbf{Үгүй.} Тогтвортой санах ой — ямар ч алдаанд өгөгдөл алдагдахгүй гэсэн үг. Хуулбарын тоог хэчнээн олон болгож, өөр өөр газар байршуулсан ч бүх хуулбарыг нэгэн зэрэг устгах үйл явдал (их хэмжээний гамшиг, программын алдаа бүх хуулбарыг гэмтээх) онолын хувьд үргэлж боломжтой. Бодит байдал дээр олон бие даасан хуулбар (RAID, өөр байршлын нөөц) ба дараалсан, шалгалттай бичих дүрмээр тогтвортой санах ойд \emph{дурын хэмжээгээр ойртуулж} болно — гэхдээ 100\% биш.

%---------------------------------------------------------------
\question{11.7}{L2 кэш: хандалтын хүлээлт 8 нс, урсгалын дамжуулах хурд 800 MB/с; үндсэн санах ой: 60 нс, 80 MB/с; хатуу диск: 15 мс, 5 MB/с; соронзон хальс: 60 с, 2 MB/с. (a) Дискний хувьд дундаж хандалтын дараа (1) 512 байт, (2) 8 KB, (3) 1 MB, (4) 16 MB урсгал дамжуулалт хийвэл үр ашигтай хурд хэд вэ? (b) Тус бүрийн ашиглалт (үр ашигтай / урсгалын хурд) хэд вэ? (c) 25\% ба түүнээс дээш ашиглалтыг хүлээн зөвшөөрөгдөхүйц гэвэл дискний хамгийн бага дамжуулалтын хэмжээ хэд вэ? (d) “Диск нь \underline{\quad} байтаас том дамжуулалтад random-access, жижигт sequential-access төхөөрөмж” гэдгийг нөхөж бич. (e) Кэш, санах ой, соронзон хальсны хувьд хүлээн зөвшөөрөгдөхүйц ашиглалтын хамгийн бага хэмжээг ол. (f) Соронзон хальс хэзээ random-access, хэзээ sequential-access төхөөрөмж вэ?}

Хүлээлт $L$, урсгалын хурд $R$, дамжуулалтын хэмжээ $x$ бол: үр ашигтай хурд $E = \dfrac{x}{L + x/R}$, ашиглалт $U = \dfrac{E}{R} = \dfrac{x}{LR + x}$. (Хэмжээг хоёртын — KB $= 2^{10}$, MB $= 2^{20}$ байт, хурдыг аравтын — MB/с $= 10^6$ байт/с гэж авав.)

\textbf{(a), (b)} Диск, $L = 15$ мс, $R = 5 \cdot 10^6$ байт/с:
\begin{center}
\begin{tabular}{lrrr}
\toprule
Хэмжээ & Дамжуулах хугацаа & Үр ашигтай хурд & Ашиглалт\\
\midrule
512 байт & 0.10 мс & 33.9 KB/с & 0.68\%\\
8 KB & 1.64 мс & 492 KB/с & 9.85\%\\
1 MB & 209.7 мс & 4.67 MB/с & 93.3\%\\
16 MB & 3355 мс & 4.98 MB/с & 99.6\%\\
\bottomrule
\end{tabular}
\end{center}

\textbf{(c)} $U \ge 0.25 \iff x \ge \dfrac{LR}{3}$. Диск: $x \ge \dfrac{0.015 \cdot 5 \cdot 10^6}{3} = \mathbf{25\,000}$ байт ($\approx 25$ KB).

\textbf{(d)} Диск нь \textbf{$\approx 25$ KB}-аас том дамжуулалтад random-access, түүнээс жижигт sequential-access төхөөрөмж (жижиг хандалтын хувьд хүлээлт давамгайлж, дараалан унших л үр ашигтай).

\textbf{(e)}
\begin{center}
\begin{tabular}{lrrr}
\toprule
Төхөөрөмж & $L$ & $R$ & Хамгийн бага $x = LR/3$\\
\midrule
L2 кэш & 8 нс & 800 MB/с & $\approx 2.1$ байт\\
Үндсэн санах ой & 60 нс & 80 MB/с & $1.6$ байт\\
Соронзон хальс & 60 с & 2 MB/с & $40 \cdot 10^6$ байт ($\approx 40$ MB)\\
\bottomrule
\end{tabular}
\end{center}
Кэш ба санах ой бараг аливаа хэмжээний (хэдэн байт) хандалтад random-access байна.

\textbf{(f)} Соронзон хальс нь $\approx 40$ MB-аас том дамжуулалтад random-access (хүлээлтийн 60 с нь дамжуулалтын хугацаатай харьцуулахад бага), түүнээс жижиг дамжуулалтад sequential-access төхөөрөмж.

%---------------------------------------------------------------
\question{11.8}{RAID level 1 зохион байгуулалт уншилтын хүсэлтэд RAID level 0 (давхардалгүй хуваалт)-оос илүү гүйцэтгэлтэй байж чадах уу? Чадвал хэрхэн?}

\textbf{Тийм.} RAID 1-д блок бүр хоёр дискэнд толин хуулбартай тул уншилтыг аль ч дискээс хийж болно: (1) хоёр өөр уншилтын хүсэлтийг хоёр дискээр \emph{зэрэг} гүйцэтгэнэ — нэг блок дээр ч гэсэн; (2) толгой нь хүссэн цилиндрт ойр байгаа (хайлт ба эргэлтийн хүлээлт бага) дискийг сонгоно. RAID 0-д блок бүр зөвхөн нэг дискэнд байдаг тул ижил дискэнд оногдсон хоёр хүсэлт дараалалд зогсож, сонголт хийх боломжгүй. (Бичилтэд RAID 1 давуу биш — хоёр хуулбарт хоёуланд нь бичнэ.)

%---------------------------------------------------------------
\question{11.9}{Хоёрдогч санах ойд HDD ашиглах гурван шалтгаан хэл.}

\begin{enumerate}[leftmargin=1.5em]
  \item \textbf{Нэг GB-ийн үнэ хямд} — NVM-ээс хэд дахин хямд тул их хэмжээний өгөгдөл, нөөц хадгалахад ашигтай.
  \item \textbf{Их багтаамж} — нэг төхөөрөмжид олон TB.
  \item \textbf{Элэгдэлгүй дахин бичигдэх} — NVM-ийн флаш нүд бичих/устгах циклийн хязгаартай (wear) бол HDD-д тийм хязгаар байхгүй; мөн дараалсан (sequential) уншилт, бичилтэд хангалттай хурдтай, удаан хугацаанд хэрэглэгдэж ирсэн найдвартай технологи.
\end{enumerate}

%---------------------------------------------------------------
\question{11.10}{Хоёрдогч санах ойд NVM төхөөрөмж ашиглах гурван шалтгаан хэл.}

\begin{enumerate}[leftmargin=1.5em]
  \item \textbf{Хурд} — хайлт, эргэлтийн хүлээлтгүй тул хандалтын хүлээлт микросекундын түвшинд, санамсаргүй (random) хандалт HDD-ээс олон зуу дахин хурдан, IOPS өндөр.
  \item \textbf{Хөдөлгөөнт хэсэггүй} — цохилт, чичиргээнд тэсвэртэй, механик эвдрэл байхгүй тул найдвартай, чимээгүй.
  \item \textbf{Бага эрчим хүч, жижиг хэмжээ} — зөөврийн компьютер, утас, хавтгай (M.2) хэлбэрт тохиромжтой; дулаан бага ялгаруулна.
\end{enumerate}

%---------------------------------------------------------------
\section{Programming Problem 11.27 (бичлэгт)}

\texttt{lab10/disk\_sched.c}: 0–4999 дугаартай 5000 цилиндртэй диск, санамсаргүй 1000 цилиндрийн хүсэлт үүсгэж FCFS, SCAN, C-SCAN-аар үйлчлээд толгойн нийт шилжилтийг хэвлэнэ. Толгойн эхний байрлал командын мөрөөс.
\begin{itemize}[leftmargin=1.5em]
  \item \textbf{FCFS} — ирсэн дарааллаар: $\sum |r_i - r_{i-1}|$.
  \item \textbf{SCAN} (цахилгаан шат) — номын жишээ шиг эхлээд 0 руу, замдаа доорх хүсэлтүүдийг үйлчилж, 0 дээр эргээд дээш хамгийн том хүсэлт хүртэл: толгой $+$ хамгийн том хүсэлт.
  \item \textbf{C-SCAN} — дээш 4999 хүртэл, тэндээс 0 руу шууд буцаж (буцах шилжилтийг нийт дүнд \emph{оруулсан}), дахин дээш толгойноос доорх хамгийн том хүсэлт хүртэл.
\end{itemize}
\begin{lstlisting}[language={}]
$ ./disk_sched 53 -c 200 -r "98 183 37 122 14 124 65 67"   # номын жишээ
FCFS:   total head movement = 640 cylinders
SCAN:   total head movement = 236 cylinders
C-SCAN: total head movement = 382 cylinders
$ ./disk_sched 2150 -s 7
Cylinders 0..4999, 1000 requests, initial head 2150
FCFS:   total head movement = 1654057 cylinders
SCAN:   total head movement = 7148 cylinders
C-SCAN: total head movement = 9996 cylinders
\end{lstlisting}
Номын §11.2-ийн жишээгээр FCFS 640, SCAN 236 нь номын хариутай таарч байна. 1000 санамсаргүй хүсэлтэд FCFS толгойг дунджаар хүсэлт бүрт $\approx 1650$ цилиндр зөөдөг бол SCAN/C-SCAN нийтдээ дискийг ердөө 1.5–2 удаа гатална — 200 гаруй дахин бага. Үр дүнг \texttt{diskcalc.py}-ийн алхам алхмын тусдаа симуляцаар шалгасан.

\end{document}
